\section{Vulnérabilité et sécurité des données}
La sécurité informatique de Cujas s'est construite sur un modèle de séparation défensive. L'infrastructure physique de la bibliothèque Cujas repose sur des serveurs fermés, un SFTP nocturne dictant et contraignant ainsi l'organisation humaine du travail, forçant les agents à travailler de manière isolée sur leurs postes plutôt que de manière collaborative. En effet, la sécurité et la conservation étaient des priorités importantes dans le \enquote{deuxième âge du numérique}. Avec l'irruption du web et la numérisation massive entraînent une multiplication du \enquote{potentiel numérique en matière de disponibilité des corpus}. Face à cet essor, les institutions patrimoniales sont aménées à réfléchir à des moyens permettant l'ouverture de leurs données tout en les protégeant. Le deuxième âge du numérique, terme employé par Emmanuelle Bermès, illustre les enjeux de cette période notamment face aux \enquote{enjeux de structuration de l'information au format numérique, d'interopérabilité et de pérennisation}. Cela permet aux institutions patrimoniales de mettre à disposition leurs contenus numérisés.

De plus, pour permettre cette diffusion du savoir, il est nécessaire d'utiliser les mêmes standards de structuration de la donnée, comme le MARC afin de rédiger les notices bibliographiques. Ces règles permettent aujourd'hui la transmission des documents numérisés entre les différentes bibliothèques numériques. Toutefois, ces concepts ont été mis en place d'une façon si stricte dans la bibliothèque Cujas qu'ils paralysent aujourd'hui le travail collaboratif des agents et figent les métadonnées. 

Cette hypervigilance de la part de la bibliothèque Cujas résulte d'un traumatisme de cyberattaque révélant ainsi une vulnérabilité locale. En 2017, la bibliothèque Cujas a subi une tentative de ransomware par cryptage. L'attaque a entraîné la \enquote{perte de fichiers, et quelques documents récupérés grâce au CINES}. Heureusement, \enquote{l'essentiel des fichiers images ont pu être sauvés (retrouvés dans des sotckages), mais le travail sur les fichiers d'une cinquantaine de livres (contrôle qualité, traitement...) ont dû être refaits}. Cela entraîne une crainte amenant à devenir davantage prudent sur les sauvegardes de la bibliothèque Cujas. Suite à cette attaque, \enquote{des sauvegardes journalières ont été mises en place de l'espace de traitement}. Cet événement illustre bien la fragilité des données \enquote{de traitement} par rapport aux données \enquote{d'archivage}. Ce traumatisme a eu pour effet de pousser l'informatique a verrouiller la chaîne de production au détriment de l'ergonomie. 

Cette sécurisation a entraîné un alourdissement de la chaîne de numérisation amenant à une saturation de leur \enquote{baie de numérisation}. En effet, la bibliothèque Cujas utilise un espace de stockage local permettant de cloisonner les différentes du processus présent dans la chaîne de numérisation. Chacun des espaces de stockage représente une tâche à effectuer. Cet espace de stockage est constitué de NAS, possédant un espace de stockage massif entre 40 et 50 téraoctets. Aujourd'hui, cet espace se retrouve saturé, encombré de répertoires multiples qui rendent la structure globale illisible. Cette saturation résulte notamment du blocage lié à l'impossibilité pour le moment des fichiers vers le CINES. La bibliothèque est donc obligé de stocker tous les paquets SIP qui doivent être transmis au CINES dans leur propre espace de stockage interne. Ce contretemps amène la bibliothèque Cujas a devoir repensé sa chaîne de numérisation afin d'intégrer par exemple des formats d'images plus léger afin de libérer de l'espace de stockage en attendant de pouvoir à nouveau envoyer leur SIP au CINES. En effet, ils envisagent probablement de passer du format TIFF au stockage au format JPEG. De nombreuses bibliothèques patrimoniales ont également été confronté à cette question. Ils ont choisi de faire ce basculement afin de libérer l'espace, tout en leur permettant de réduire leur empreinte carbonne liée à leur projet de numérisation.

Le poids des fichiers stockés dans les espaces de stockages n'est pas le seul problème rencontré par Cujas, ils doivent également faire face à la chaîne de numérisation. Le travail des experts métiers de la bibliothèque se trouve segmenté à tel point qu'afin d'alléger légèrement les démarches de travail. Des moulinettes sont été crées afin de permettre une forme d'automatisation dans le travail des ouvrages numérisés. Ces moulinettes aident à combler des lacunes identifiées par le personnel de Cujas qui a été mis en place par le responsable DIT les soulageant par exemple dans la production des notices numériques. Cependant, le moindre changement peut entraîner une casse des \enquote{moulinettes} car correspondant à un schéma particulier de travail des équipes. Elles sont donc difficilement adaptable ou demandent une profonde réflexion afin d'améliorer au mieux la production. 

En réaction à l'attaque de 2017, le responsable du DIT a énormément sécurisés les accès des équipes allant jusqu'à différencier les droits d'accès de chacun concernant les différentes étapes de production. Personne n'a accès à tout les paramètres hormis lui. La publication des données ne peut plus se faire de façon simultanée car celles-ci passent désormais par un outil SFTP sécurisé tels que WinSCP avec un traitement par lot. Ces traitements se réalisent uniquement la nuit car selon le responsable du DIT c'est le seul moment où les transferts peuvent s'effectuer sans trop impactés le reste du travail. En journée, cela n'est pas possible à cause de l'occupation par tous le personnel des espaces amenant des traitements plus long que lorsqu'ils sont effectués de nuit.

L'espace sur lesquels le personnel de Cujas travaille est la baie de numérisation et celle-ci est accessible uniquement aux personnes possédant les codes d'accès sur des postes prédéfinis afin d'effectuer une tâche précise. La structure de travail étant instable car elle a été crée au début comme un espace de stockage ne devant pas autant de connection que ce qu'elle a maintenant amenant à une vigilance des équipes lorsqu'ils se connectent dessus. Il leur est interdit d'utiliser un moteur de recherche à l'intérieur afin de ne pas la fragiliser et de protéger les données qui y sont traitées. Par ailleurs, il n'est pas permis aux équipes de s'y connecter à distance que ce soit dans un autre lieu de la bibliothèque ou chez eux. La baie de numérisation n'ayant pas été suffisamment balisé demande une prudence dans l'exécution de leur tâche afin de ne pas compromettre les données. 
Cette surprotection des équipes de Cujas via leur chaîne de numérisation amène à réfléchir à la notion même de sécurisation autour de la donnée. En effet, à une époque où l'on pense à ouvrir celle-ci permettant à d'autres institutions patrimoniales de pouvoir l'utiliser pour leur propre collection. Cela amène à s'interroger sur les conditions de sécurités et de préservation nécessaire sans que cela ne bloque complètement cet échange de flux de données qui voit le jour avec l'arrivée du Web sémantique. Bien sûr, la sécurité et la préservation sont des conditions nécessaires et importantes pour les bibliothèques mais ne doivent pas pour autant entraîner un isolement des données de l'institution à cause d'un modèle de sécurité trop stricte. Ce modèle de sécurité utilisé par Cujas est typique des silos du \enquote{deuxième âge}, créant une contradiction avec les méthodes de travail modernes empêchant une autonomie relative de la donnée. 

Cette hypervigilance justifiée à bien des égards entraîne un isolement de la bibliothèque Cujas par rapport aux autres bibliothèques numériques. La rigidité de leur condition de travail entraîne comme conséquence l'impossibilité de faire évoluer leur système pour répondre aux attentes d'aujourd'hui. Les exigences de maintenant sont davantages flexible avec le Web de données que dans les années 1990 et 2010 empêchant la bibliothèque d'améliorer l'ergonomie de son interface.